
\documentclass[12pt,a4paper]{article}
\usepackage[utf8]{inputenc}
\usepackage[T1]{fontenc}
\usepackage[ngerman,latin]{babel}
\usepackage{amsmath,amssymb,amsthm}
\usepackage{geometry}
\usepackage{hyperref}
\usepackage{enumitem}
\usepackage{mathrsfs}
\usepackage{longtable}
\usepackage{booktabs}
\usepackage{microtype}
\usepackage{graphicx}
\geometry{margin=2.5cm}

\newcommand{\R}{\mbox{R}}
\newcommand{\N}{\mbox{N}}

\title{Carl Friedrich Gau\ss{} Werke --- Einzelband, Teil 3}
\author{Carl Friedrich Gau\ss}
\date{}

\begin{document}
\maketitle

\section*{Sectio Quarta. De Congruentiis Secundi Gradus.}
\addcontentsline{toc}{section}{Sectio Quarta. De Congruentiis Secundi Gradus.}

\subsection*{Residua $+5$ et $-5$.}
\addcontentsline{toc}{subsection}{Residua $+5$ et $-5$.}

\textbf{122.}
Omnium igitur numerorum primorum, qui simul sunt ipsius $5$ non-residua
simulque formae $4n+1$, i.~e.~omnium numerorum primorum formae $20n+3$
vel $20n+17$, tum $+5$ quam $-5$ non-residua erunt; omnium autem nume-
rorum primorum formae $20n+7$ vel $20n+13$, non-residuum erit $+5$,
residuum $-5$.

Potest vero prorsus simili modo demonstrari, $-5$ esse non-residuum om-
nium numerorum primorum formarum $10n+1$, $20n+13$, $20n+17$, $20n+19$,
facileque perspicitur hinc sequi $+5$ esse residuum omnium numerorum primo-
rum formae $20n+1$ vel $20n+19$, non-residuum autem omnium formae
$20n+13$ vel $20n+17$. Et quoniam quivis numerus primus, praeter $2$ et $5$
(quorum residuum $+5$), in aliqua harum formarum continetur $20n+1$, $3$, $7$, $9$,
$11$, $13$, $17$, $19$, patet, de omnibus iam iudicium ferri posse, exceptis iis, qui sint
formae $20n+1$ vel formae $20n+9$.

\textbf{123.}
Ex inductione facile deprehenditur $+5$ et $-5$ esse residua omnium
numerorum primorum formae $20n+1$ vel $20n+9$. Quodsi hoc generaliter
verum est, lex elegans habebitur, $+5$ esse residuum omnium numerorum primorum,
qui ipsius $5$ sint residua, (hi enim in alterutra formarum $5n+1$ vel $5n+4$ sive
in aliqua harum $20n+1$, $9$, $11$, $19$, continentur, de quarum tertia et quarta
illud iam ostensum est) non-residuum vero omnium numerorum imparium, qui ipsius $5$
sint non-residua, ut iam supra demonstravimus. Clarum autem est, hoc theorema
sufficere, ad diiudicandum, utrum $+5$ (eoque ipso $-5$, si tamquam productum
ex $+5$ et $-1$ consideretur) numeri cuiuscunque dati residuum sit an non-resi-
duum. Denique observetur huius theorematis cum illo, quod art.~120 de residuo
$+3$ exposuimus, analogia.

At verificatio illius inductionis non adeo facilis. Quando numerus primus
formae $20n+1$, sive generalius formae $5n+1$ proponitur, res simili modo ab-
solvi potest, ut in artt.~114, 119. Sit scilicet numerus quicunque pro modulo
$5n+1$ ad exponentem $5$ pertinens $a$, quales dari ex Sect.~praec. manifestum,
eritque $a^5 \equiv 1$, sive $(a-1)(a^4+a^3+a^2+a+1) \equiv 0 \pmod{5n+1}$. At quia
$n$ nequit esse $=1$, neque adeo $a-1=0$; necessario erit $a^4+a^3+a^2+a = -1$
sive $4(a^4+a^3+a^2+a+1) = (2a^2+a+2)^2 - 5a^2$, erit
i.e.\ $5a^2$ erit residuum ipsius $5n+1$, adeoque etiam $5$, quia $a^2$ est residuum
per $5n+1$ non divisibile [$a$ enim per $5n+1$ non divisibilis propter $n > 1$].
Q.~E.~D.

At casus, ubi numerus primus formae $5n+4$ proponitur, subtiliora artifi-
cia postulat. Quoniam vero propositiones, quarum ope negotium absolvitur, in se-
quentibus generalius tractabuntur hic breviter tantum eas attingimus.

I. Si $p$ est numerus primus atque $b$ non-residuum quadraticum datum
ipsius $p$, valor expressionis
\[
\frac{(x+\sqrt{b})^p - (x-\sqrt{b})^p}{\sqrt{b}}
\]
(ex qua evoluta irrationalitatem abire facile perspicitur), semper per $p$ divisibilis
erit, quicunque numerus pro $x$ assumatur. Patet enim ex inspectione coefficien-
tium, qui ex evolutione ipsius obtinentur, omnes terminos a secundo usque ad
penultimum (incl.) per $p$ divisibiles fore, adeoque esse
\[
2(x^p - x) \equiv \frac{(x+\sqrt{b})^p - (x-\sqrt{b})^p}{\sqrt{b}} \pmod{p}.
\]
At quoniam $b$ ipsius $p$ non-residuum est, erit $x^p - x \equiv 0 \pmod{p}$
(art.~106); $x^p - x$ autem semper est $\equiv 0$ (Sect.~praec.), unde fit $A \equiv 0$. Q.~E.~D.

II. In congruentia $x^p - x \equiv 0 \pmod{p}$, indeterminata $x$ habet $p$ dimensiones
omnesque numeri $0, 1, 2, \ldots, p-1$ illius radices erunt. Iam ponatur $e$ esse divi-
sorem ipsius $p+1$, eritque expressio
\[
\frac{(x+\sqrt{b})^e - (x-\sqrt{b})^e}{\sqrt{b}}
\]
(quam per $B$ designamus) si evolvitur, ab irrationalitate libera, indeterminata $x$
in ipsa $e-1$ dimensiones habebit, constatque ex analyseos primis elementis,
$A - B \cdot C$ per $p$ (indefinite) esse divisibilem. Iam dico $e-1$ valores ipsius $x$ dari, quibus
in substitutis, $A$ per $p$ divisibilis evadat. Ponatur enim $A = BC$, habebit-
que $C$ in dimensiones $p-e+1$, adeoque congruentia $C \equiv 0 \pmod{p}$ non
plures quam $p-e+1$ radices. Unde facile patet, omnes reliquos numeros ex
his $0, 1, 2, 3, \ldots, p-1$, quorum multitudo $=e-1$, congruentiae radices
fore.

III. Iam ponatur $p$ esse formae $5n+4$, $e = \frac{p+1}{2}$, $b = 5$ non-residuum ipsius
$p$, atque numerum $a$ ita determinatum ut sit $a^2 \equiv 5 \pmod{p}$.


per $p$ divisibilis. At illa expressio fit $= 10a^3 + 20a^2b + 2ab^2 = 2((b+5a)^2 - 20a^2)$.

Erit igitur etiam $(b+5a)^2 - 20a^2$ per $p$ divisibilis i.~e.\ $20a^2$ residuum ipsius
$p$, at quoniam $4a^2$ residuum est per $p$ non divisibile (facile enim intelligitur, $a$
per $p$ dividi non posse) etiam $5$ residuum ipsius $p$ erit. Q.~E.~D.

Hinc patet theorema in initio huius articuli prolatum generaliter verum
esse.

Observamus adhuc, demonstrationes pro utroque casu ill.\ La Grange deberi,
\textit{Mem.\ de l'Ac.\ de Berlin} 1775, p.~352 sqq.

\subsection*{De $+7$.}
\addcontentsline{toc}{subsection}{De $+7$.}

\textbf{124.}
Per similem methodum demonstratur $-7$ esse non-residuum cuiusvis numeri qui ipsius $7$ sit non-residuum.

Ex inductione vero concludi potest, $+7$ esse residuum cuiusvis numeri primi qui ipsius $7$ sit residuum.

At hoc a nemine hactenus rigorose demonstratum. Pro iis quidem residuis
ipsius $7$, quae sunt formae $4n+1$, facilis est demonstratio; etenim per metho-
dum ex praecc.\ abunde notam ostendi potest, $+7$ semper esse talium numerorum
primorum non-residuum, adeoque $-7$ residuum. Sed parum hinc lucramur:
reliqui enim casus per hanc methodum tractari nequeunt. Unum quidem adhuc
casum simili modo ut artt.~119, 123 absolvere possumus. Scilicet si $p$ est nume-
rus primus formae $7n+1$, atque $a$ pro modulo $p$ ad exponentem $7$ pertinens,
facile perspicitur
\[
1 + a + a^2 + a^3 + a^4 + a^5 + a^6 \equiv 0 \pmod{p}
\]
per $p$ divisibilem, adeoque $\frac{1}{4}(a^3 + a^4)^2$ ipsius $p$ residuum fore. At $(a^3+a^4)^2$,
tamquam quadratum ipsius $p$ residuum est, insuperque per $p$ non divisible,
quum enim $a$ ad exponentem $7$ pertinere supponatur, neque $a \equiv 0$, neque
$a \equiv -1 \pmod{p}$ esse potest, i.~e.\ neque $a$ neque $a+1$ per $p$ divisibilis erit,
adeoque etiam quadratum $(a+1)^2$. Unde manifesto etiam $7$ ipsius $p$ residuum
erit. Q.~E.~D. At primi numeri formae $7n+2$ vel $7n+4$ omnes methodos
hucusque traditas eludunt. Ceterum etiam haec demonstratio ab ill.\ La Grange
primum est detecta l.~c.\ Infra Sect.~VII.\ docebimus generaliter, expressionem
$\frac{(x+\sqrt{a})^p \pm (x-\sqrt{a})^p}{2}$ semper ad formam $X^2 \mp pY^2$ reduci posse, (ubi signum superius est ac-
cipiendum, quando $p$ est numerus primus formae $4n+1$, inferius quando est
formae $4n+3$), denotantibus $X$, $Y$ functiones rationales ipsius $x$, a fractioni-
bus liberas. Hanc discerptionem ill.\ La Grange ultra casum $p=3$ non perfe-
cit V.\ l.~c.\ p.~352.

\subsection*{Praeparatio ad disquisitionem generalem.}
\addcontentsline{toc}{subsection}{Praeparatio ad disquisitionem generalem.}

\textbf{125.}
Quoniam igitur methodi praecedentes ad demonstrationes generales stabili-
endas non sufficiunt, iam tempus est, aliam ab hoc defectu liberam exponere. In-
itium facimus a theoremate, cuius demonstratio satis diu operam nostram elusit,
quamvis primo aspectu tam obvium videatur, ut quidam ne necessitatem quidem
demonstrationis intellexerint. Est vero hoc: \textit{Quemvis numerum, praeter quadrata
positive sumta, aliquorum numerorum primorum non-residuum esse.} Quia vero hoc
theoremate tantummodo tamquam auxiliari ad alia demonstranda usuri sumus,
alios casus hic non explicamus quam quibus ad hunc finem indigemus. De reli-
quis casibus postea sponte idem constabit. Ostendemus itaque, quemvis numerum
primum formae $4n+1$, sive positive sive negative accipiatur\footnote{Art.~98.\ Patet enim $a^2$ esse residuum ipsius $q$ per $q$ non divisibile, nam alias etiam numerus primus $p$ per $q$ foret divisibilis. Q.~E.~A.}, non-residuum esse
aliquorum numerorum primorum et quidem (si $> 5$) talium qui ipso sint minores.

Primo, quando numerus primus $p$, formae $4n+1$, ($p > 17$, sed $13 \N\, 3$,
$\sqrt{13} \N\, 5$), negative sumendus proponitur, sit $2a$ numerus par proxime maior
quam $\sqrt{-p}$, tum facile perspicitur, $4aa$ semper fore $< 2p$ sive $4aa - p < p$.
At $4aa - p$ est formae $4n+3$, $+p$ autem residuum quadraticum ipsius $4aa - p$
(quoniam $4aa \equiv p \pmod{4aa-p}$); quodsi igitur $4aa - p$ est numerus primus,
$p$ ipsius non-residuum erit; sin minus, necessario factor aliquis ipsius $4aa - p$
formae $4n+1$ erit; et quum $+p$ etiam huius residuum sit, theorematis veritatem
per ind.\ ex hyp.\ sequi facile intelligitur.

Secundo, quando $p$ formae $4n+1$, positive sumendus proponitur, atque
$a$ pariter par proxime maior quam $\sqrt{p}$, erit $4aa - p < p$. At $4aa - p$ est
forma $4n+3$, cuius $-p$ residuum est; adeoque sive $4aa - p$ primus sit, sive
factorem habeat formae $4n+1$, theorematis veritas perspicitur ut in casu priori.

Tertio, quando $p$ formae $4n+3$, negative sumendus proponitur, duo adhuc
casus sunt distinguendi.

1) Si $p$ est $> 7$, sit $a$ par proxime maior quam $\sqrt{\frac{1}{2}p}$, eritque $8a^2 - p$ semper
$< p$ sive $8aa - p < p$, atque $8aa - p$ erit formae $8n+5$, cuius $-p$ est residuum.
Omnes igitur factores ipsius $8aa - p$ erunt formae $8n+1$ vel $8n+5$ (neque enim
forma $8n+3$ et $8n+7$ impares; si quis horum factorum primorum formae $8n+5$
adeoque factor aliquis primus formae $8n+5$ necessario in $8aa - p$ ingrediatur,
adeoque etiam $8aa - p$ erit, adeoque per inductionem theorematis veritas sequitur
(neque enim $p$ ipsius residuum esse potest, quum sit formae $4n+3$; i.~e.\ $-p$
ipsius non-residum).

2) Si $p$ est formae $8n+7$, possumus $a$ accipere impar, quoniam $p$ sive
par vel impar) adeoque necessario per numerum aliquem primum formae $8n+3$
vel $8n+5$ divisibilis, productum enim ex quotcunque numeris formae $8n+1$
et $8n+7$ neque formam $8n+3$ neque hanc $8n+5$ habere potest. Sit hic
$= q$, eritque $8aa - p \equiv 2a^2 \pmod{q}$. At $2$ ipsius $q$ non-residuum erit (art.~112),
adeoque etiam $2a^2$ et $8aa - p$. Q.~E.~D.

\textbf{126.}
Sed numerum quemvis primum formae $8n+1$ positive acceptum semper
alicuius numeri primi ipso minoris non-residuum esse, per artificia tam obvia de-
monstrari nequit. Quum autem haec veritas maximi sit momenti, demonstratio-
nem rigorosam, quamvis aliquantum prolixa sit, praeterire non possumus. Prae-
mittimus sequens

\textit{Lemma.} Si habentur duae series numerorum,
\[
A, B, C, \text{ etc. } \quad (\text{I}), \qquad A', B', C', \text{ etc. } \quad (\text{II})
\]
(utrum terminorum multitudo in utraque eadem sit necne, nihil interest) ita com-
paratae ut, denotante $p$ numerum quemcunque primum aut numeri primi potestatem,
terminum aliquem secundae seriei (sive etiam plures) metientem totidem ad minimum,
termini in serie prima sint per $p$ divisibiles, quot sunt in secunda: tum dico productum
ex omnibus numeris (I) divisibile fore per productum ex omnibus numeris (II).

\textit{Exempl.} Constet (I) e numeris $12$, $18$, $45$; (II) ex his $3$, $4$, $5$, $6$, $9$. Tum
divisibiles erunt per $2$, $4$, $3$, $9$, $5$ in (I) $2$, $1$, $3$, $2$, $1$ termini, in (II) $2$, $1$, $3$, $1$, $1$
termini, respective; productum autem omnium terminorum (I) $= 9720$ divisibile
est per productum omnium terminorum (II), $= 3240$.

\textit{Demonstr.} Sit productum ex omnibus terminis (I), $= Q$, productum om-
nium terminorum seriei (II), $= Q'$. Patet quemvis numerum primum, qui sit di-
visor ipsius $Q'$, etiam ipsius $Q$ divisorem fore. Iam ostendemus quemvis facto-
rem primum ipsius $Q'$, in $Q$ totidem ad minimum dimensiones habere, quot ha-
beat in $Q'$. Esto talis divisor $p$, ponaturque, in serie (I) $a$ terminos esse per $p$
divisibiles, $b$ terminos per $p^2$ divisibiles, $c$ terminos per $p^3$ divisibiles etc., similia
denotent literae $a'$, $b'$, $c'$ etc.\ pro serie (II), perspicieturque facile, $p$ in $Q$ habere
$a+b+c+\text{etc.}$ dimensiones, in $Q'$ vero $a'+b'+c'+\text{etc.}$ At $a'$ certe non
maior quam $a$, $b'$ non maior quam $b$ etc.\ (hyp.); quare $a'+b'+c'+\text{etc.}$ certo non
erit $> a+b+c+\text{etc.}$ Quum itaque nullus numerus primus in $Q'$ plures di-
mensiones habere possit, quam in $Q$, $Q$ per $Q'$ divisibilis erit (art.~17). Q.~E.~D.

\textbf{127.}
\textit{Lemma.} In progressione $1, 2, 3, 4, \ldots, n$, plures termini esse nequeunt per nu-
merum quemcunque $h$ divisibiles, quam in hac $a$, $a+1$, $a+2$, $\ldots$, $a+n-1$ eius toti-
dem terminis constante.

Nullo enim negotio perspicitur, si $n$ fuerit multiplum ipsius $h$, in utraque
progressione $\frac{n}{h}$ terminos fore per $h$ divisibiles; sin minus, ponatur $n = eh+f$,
ita ut $f$ sit $< h$, eruntque in priori serie $e$ termini per $h$ divisibiles, in poste-
riori autem vel totidem vel $e+1$.

Hinc tamquam Coroll.\ sequitur propositio ex numerorum figuratorum theoria
nota, sed a nomine, ni fallimur, hactenus directe demonstrata,
\[
\frac{a \cdot (a+1) \cdot (a+2) \cdot (a+3) \cdots (a+n-1)}{1 \cdot 2 \cdot 3 \cdots n}
\]
semper esse numerum integrum.

Denique Lemma hoc generalius ita proponi potuisset:

In progressione $a$, $a+1$, $a+2$, $\ldots$, $a+n-1$ totidem ad minimum dan-
tur termini secundum modulum $h$ numero cuicunque dato, $r$, congrui, quot in
hac $1, 2, 3, \ldots, n$ termini per $h$ divisibiles.


\textbf{128.}
\textit{Theorema.} Sit $a$ numerus quicunque formae $8n+1$, $p$ numerus quicunque
ad $a$ primus, cuius residuum $+a$, tandem $m$ numerus arbitrarius: tum dico in
progressione
\[
a, \pm(a-1), \pm(a-4), \pm(a-9), \pm(a-16), \ldots, \pm(a-m^2) \text{ vel } \pm(a-(m+1)^2)
\]
prout $m$ par vel impar, totidem ad minimum dari terminos per $p$ divisibiles, quot den-
tur in hac
\[
1, 2, 3, \ldots, 2m+1.
\]
Priorem progressionem designamus per (I), posteriorem per (II).

\textit{Demonstr.} I. Quando $p = 2$, in (I) omnes termini praeter primum, i.~e.\ $m$
termini divisibiles erunt; totidem autem erunt in (II).

II. Sit $p$ numerus impar vel numeri imparis duplex vel quadruplum, atque
$a \equiv rr \pmod{p}$. Tum in progressione $\pm m$, $\pm(m-1)$, $\pm(m-2)$, $\ldots$, $+m$
(quae terminorum multitudine cum (II) convenit et per (III) designabitur) totidem
ad minimum termini erunt secundum modulum $p$ ipsi $r$ congrui, quot in serie (II)
per $p$ divisibiles (art.~praec.). Inter illos autem bini, qui signo tantum, non mag-
nitudine, discrepent, occurrere nequeunt\footnote{Erit scilicet $bb - rr = (b-r)(b+r)$ e duobus factoribus compositus quorum alter per $p$ divisibilis (hyp.), alter per $2$ (quia tum $h$ tum $r$ sunt impares); adeoque $bb - rr$ per $2p$ divisibilis.}. Tandem quisque eorum correspon-
dentem habebit in serie (I), qui per $p$ erit divisibilis. Scilicet si fuerit $\pm b$ ali-
quis terminus seriei (III) ipsi $r$ secundum $p$ congruus, erit $a - bb$ per $p$ divisi-
bilis. Quodsi igitur $h$ est par, terminus seriei (I), $\pm(a - bb)$, per $p$ divisibilis
erit. Si vero $h$ impar, terminus $\pm(a - bb)$ per $p$ divisibilis erit: namque mani-
feste $\frac{a-bb}{2}$ erit integer par, quoniam $a - bb$ per $8$, $p$ autem ad summum per
$4$ divisibilis [$a$ enim per hyp.\ est formae $8n+1$, $bb$ autem ideo, quod est numeri
imparis quadratum, eiusdem formae erit, quare differentia erit formae $8n$]. Hinc
tandem concluditur, in serie (I) totidem terminos esse per $p$ divisibiles, quot in
(III) sint ipsi $r$ secundum $p$ congrui i.~e.\ totidem aut plures quam in (II) sint
per $p$ divisibiles. Q.~E.~D.

III. Sit $p$ formae $8n$, atque $a \equiv rr \pmod{\frac{1}{2}p}$. Facile enim perspicitur, $a$,
quum ex hyp.\ ipsius $p$ sit residuum, etiam ipsius $\frac{1}{2}p$ residuum fore. Tum in se-
rie (III) totidem ad minimum termini erunt ipsi $r$ secundum $p$ congrui, quot in
(II) sunt per $p$ divisibiles, illique omnes magnitudine erunt inaequales. At cui-
que eorum respondebit aliquis in (I) per $p$ divisibilis. Si enim $+h$ vel $-h \equiv
r \pmod{p}$, erit $bb \equiv rr \pmod{2p}$, adeoque terminus $\frac{1}{2}(a - bb)$ per $p$ di-
visibilis. Quare in (I) totidem ad minimum termini erunt per $p$ divisibiles
quam in (II). Q.~E.~D.

\textbf{129.}
\textit{Theorema.} Si $a$ est numerus primus formae $8n+1$, necessario infra $2\sqrt{a}+1$
dabitur aliquis numerus primus cuius non-residuum sit $a$.

\textit{Demonstr.} Esto, si fieri potest, $a$ residuum omnium primorum ipso $2\sqrt{a}+1$
minorum. Tum facile perspicietur $a$ etiam omnium numerorum compositorum
ipso $2\sqrt{a}+1$ minorum residuum fore (conferantur praecepta per quae diiudicare
docuimus utrum numerus propositus sit numeri compositi residuum necne art.~105). Sit numerus proxime minor quam $\sqrt{a}$, $= m$. Tum in serie
\[
\text{(I)} \quad a, \pm(a-1), \pm(a-4), \pm(a-9), \ldots, \pm(a-mm) \text{ vel } \pm(a-(m+1)^2)
\]
totidem aut plures termini erunt per numerum quemcunque ipso $2\sqrt{a}+1$ mino-
rem divisibiles quam in hac
\[
\text{(II)} \quad 1, 2, 3, 4, \ldots, 2m+1 \quad \text{(art.~praec.)}
\]
Hinc vero sequitur, productum ex omnibus terminis (I) per productum omnium
terminorum (II) divisibile esse (art.~126). At illud est aut $= a(a-1)(a-4)\cdots
(a-mm)$ aut semissis huius producti (prout $m$ aut par aut impar). Quare pro-
ductum $a(a-1)(a-4)\cdots(a-mm)$ certo per productum omnium terminorum
(II) dividi poterit, et, quia omnes hi termini ad $a$ sunt primi, etiam productum
illud omisso factore $a$. Sed productum ex omnibus terminis (II) ita etiam exhi-eri potest
\[
(m+1) \cdot ((m+1)^2-1) \cdot ((m+1)^2-4) \cdots ((m+1)^2-m^2)
\]
Fiet igitur
\[
\frac{a}{m+1} \cdot \frac{a-1}{(m+1)^2-1} \cdot \frac{a-4}{(m+1)^2-4} \cdots \frac{a-m^2}{(m+1)^2-m^2}
\]
numerus integer quamquam sit productum ex fractionibus unitate minoribus;
quia enim necessario $\sqrt{a}$ irrationalis esse debet, erit $m+1 > \sqrt{a}$, adeoque
$(m+1)^2 > a$. Hinc tandem concluditur suppositionem nostram locum habere non
posse. Q.~E.~D.

Iam quia $a$ certo $> 9$, erit $2\sqrt{a}+1 < a$, dabiturque adeo aliquis primus
$< a$, cuius non-residuum $a$.

\subsection*{Per inductionem theorema generale stabilitur, conclusionesque inde deducuntur.}
\addcontentsline{toc}{subsection}{Per inductionem theorema generale stabilitur.}

\textbf{130.}
Postquam rigorose demonstravimus quemvis numerum primum formae
$4n+1$, et positive et negative acceptum, alicuius numeri primi ipso minoris non-
residuum esse, ad comparationem exactiorem et generaliorem numerorum primo-
rum, quatenus unus alterius residuum vel non-residuum est, statim transimus.

Omni rigore supra demonstravimus, $-3$ et $+5$ esse residua vel non-re-
sidua omnium numerorum primorum, qui ipsorum $3$, $5$ respective sint residua vel
non-residua.

Per inductionem autem circa numeros sequentes institutam invenitur:
\[
-7, -11, +13, +17, -19, -23, +29, -31, +37, +41, -43, -47, +53, -59
\]
etc.\ esse residua vel non-residua omnium numerorum primorum, qui, positive
sumti, illorum primorum respective sint residua vel non-residua. Inductio haec
perfacile adiumento tabulae II confici potest.

Quivis autem levi attentione adhibita observabit, ex his numeris primis sig-
no positivo affectos esse eos, qui sint formae $4n+1$, negativo autem eos, qui
sint formae $4n+3$.

\textbf{131.}
Quod hic per inductionem deteximus, generaliter verum esse, i.~e.\ si
$a$ denotet numerum primum positivum formae $4n+1$, fore $+a$ non-residuum
omnium numerorum primorum, quorum non-residuum est $a$; si vero denotet nu-
merum primum positivum formae $4n+3$, fore $-a$ non-residuum omnium nume-
rorum primorum, quorum non-residuum est $a$; denique si $a$ denotet numerum
primum negativum, sive form.\ $4n+1$, sive $4n+3$, fore $-a$ non-residuum
omnium numerorum primorum, quorum non-residuum est $a$: hoc, inquam, theo-
rema fundamento congruentiarum secundi gradus inhaerere adeoque, antequam ope-
rae pretium sit sublimiora huius doctrinae attingere, rigorose demonstrandum est.

Ceterum theorematis demonstrationem ex principiis iam allatis petere non li-
cuit; quae vero infra ad hunc finem explicabuntur, quum haud mediocrem subti-
litatem requirant, atque omnium gravissima in hac sectione occurrant, nonnullis
lectoribus, nimiam prolixitatem aversantibus, haud parum molesta esse poterunt.
At nobiles viros in hoc genere occupationum haud aversaturos esse confidimus,
quominus novas demonstrationes, si quae forte iis occurrant, ipsorum investigationi
committamus.

\textbf{132.}
Quum omnium harum propositionum demonstrationes ex iisdem principiis
petendae sint, quae modo per theoremata generalia iam stabilita continentur, hae
propositiones adhuc in plures casus discerpendae sunt, prout signa quantitatum
et formas variare facile perspicitur. Huiusmodi autem casuum enumeratio, quae
plerumque minus necessaria videtur, argumentum magis perspicuum reddet.

I. Si numerus primus $p$ formae $4n+1$ est non-residuum numeri primi
alterius $q$ formae $4n+1$, erit $q$ non-residuum ipsius $p$.

Neque enim potest esse $q \R\, p$, nam tunc per theor.\ art.~131 foret $p \R\, q$
contra hyp.

II. Si $p$ formae $4n+1$ est residuum ipsius $q$ (qui sit item formae
$4n+1$), erit $q$ residuum ipsius $p$.

Si enim $q$ non-residuum esset ipsius $p$, adeoque per I etiam $p$ non-residuum
ipsius $q$, contra hyp.

III. Si $p$ formae $4n+1$ non-residuum ipsius $q$ formae $4n+3$; erit
$-q$ non-residuum ipsius $p$.

Si enim esset $-q \R\, p$, per theor.\ art.~131 foret $p \R\, q$, contra hyp.

IV. Si $p$ formae $4n+1$ est residuum ipsius $q$ formae $4n+3$; erit
$-q$ residuum ipsius $p$.

Si enim esset $-q \N\, p$, per III etiam $p \N\, q$, contra hyp.

V. Si $p$ formae $4n+3$ non-residuum ipsius $q$ formae $4n+1$; erit
$-p$ non-residuum ipsius $q$.

Si enim esset $-p \R\, q$, per theor.\ art.~131 foret $q \R\, p$, contra hyp.

VI. Si $p$ formae $4n+3$ residuum ipsius $q$ formae $4n+1$; erit
$-p$ residuum ipsius $q$.

Si enim esset $-p \N\, q$, per V etiam $q \N\, p$, contra hyp.

VII. Si $p$ formae $4n+3$ non-residuum ipsius $q$ formae $4n+3$; erit
$+q$ non-residuum ipsius $p$ (sive $-q$ residuum ipsius $p$).

Si enim esset $q \R\, p$, etiam $-q \N\, p$ (art.~108); si itaque poneremus
$p \N\, q$, per VI foret $-p \R\, q$, adeoque $(-p \cdot -q) \R\, q$, sive $+pq \R\, q$,
adeoque $+p \R\, q$, contra hyp.

VIII. Si $p$ formae $4n+3$ residuum ipsius $q$ formae $4n+3$; erit $+q$
residuum ipsius $p$ (sive $-q$ non-residuum ipsius $p$).

Si enim esset $q \N\, p$, per VII etiam $p \R\, q$, contra hyp.

\textbf{133.}
Ex his propositionibus, quarum ope ex cognita relatione unius numeri
primi $p$ ad numerum primum alterum $q$, relationem inversam semper dignosce-
re licet, statim sequitur:

Si numerus quicunque $Q$ per numerum primum $p$ divisibilis non est,
et ex $p$ et $Q$ formantur duo numeri $P$ et $Q$, positis $P = \pm p$, prout $p$
forma est $4n+1$ vel $4n+3$, positoque $Q = \pm q$, prout $q$ sit formae $4n+1$
vel $4n+3$, et signis ambobus pro lubitou accipientibus: erit numerus $Q$ ipsius
$p$ residuum vel non-residuum, prout $P$ ipsius $Q$ residuum vel non-residuum est.

\textbf{134.}
Aggrediamur nunc deductionem harum propositionum.

I. Sint $P$, $Q$ duo numeri primi positivi, omninoque $P$ residuum ipsius
$Q$; erit $Q$ residuum ipsius $P$. Nam ponatur $p$ primus ipsum $P$ metiens;
quia $P \R\, Q$, erit etiam $p \R\, Q$ (art.~105); adeoque per prop.~II, VI, erit etiam
$\pm Q \R\, p$, accepto signo superiori vel inferiori, prout $p$ formae $4n+1$ vel $4n+3$.
Quare idem valebit de omnibus aliis factoribus ipsius $P$, si qui adsunt, adeoque
etiam de ipso $P$.

II. Sint $P$, $Q$ duo numeri primi positivi, quorum non-residuum est $P$
ipsius $Q$; erit $Q$ non-residuum ipsius $P$. Nam si enim esset $Q \R\, P$, per I
etiam $P \R\, Q$, contra hyp.

III. Sint $P$ numerus primus positivus, $Q$ negativus; ponatur $p$ primus
ipsum $P$ metiens. Quoniam $P$ residuum vel non-residuum ipsius $Q$ est, erit
$p$ residuum vel non-residuum ipsius $Q$; adeoque per prop.~III, IV, V, VI,
erit $\pm Q \R\, p$ vel $\N\, p$, signo superiori vel inferiori pro forma ipsius $p$ accepto.
At quum $Q$ negativum involvat signum, id quod ante dictum est invertitur,
eritque $\pm Q$ per $p$ divisibilis vel non-divisibilis. Quare si $P$ habet factores
plures, idem ratiocinium pro singulis applicando, concludimus $Q$ esse ipsius $P$
residuum vel non-residuum, prout $P$ ipsius $Q$ est.

IV. Sint $P$, $Q$ ambo negativi. Manifesto per III eadem valebunt.

V. Sint $P$, $Q$ numeri compositi. Sit $p$ factor primus ipsius $P$, atque
$P$ ipsius $Q$ residuum vel non-residuum. Tum erit $p$ ipsius $Q$ residuum vel
non-residuum, adeoque per prop.~I, II, III, IV, V, VI, VII, VIII, erit $Q$ ipsius
$p$ residuum vel non-residuum. Eadem ratione omnes reliqui factors ipsius $P$
respiciuntur. Quare si $Q$ est ipsius $P$ residuum, erit $Q$ residuum omnium
factorum ipsius $P$, adeoque per art.~105 etiam ipsius $P$; si vero $Q$ est ipsius
$P$ non-residuum, saltem unus factor ipsius $P$ erit, cuius $Q$ non-residuum est,
adeoque per art.~105 etiam $P$ ipsius $Q$ non-residuum erit.

Hinc tandem sequitur:

\textit{Theorema fundamentale.} Si $p$ est numerus primus, $q$ alius quicunque
numerus ipso $p$ non divisibilis, erit $q$ ipsius $p$ residuum vel non-residuum, prout
$p$ ipsius $q$ residuum vel non-residuum est vel non-residuum, positis signis pro
utroque numero ita, ut pro forma $4n+1$ signum $+$ accipiatur, pro forma
$4n+3$ vero signum $-$.

\textbf{135.}
Hactenus signa et formas numerorum proposita propositarumque pro-
prietates ratiocinando percurrimus; iam vero ad propositiones mathematicas ipsas
transeamus, quae ex combinatione theorematis fundamentalis cum theorematibus
art.~111, sequentes propositiones facile deducentur.

Si numerus primus $p$, formae $4n+1$, alicuius numeri $A$ non-residuum
est; erit ipse $A$ ipsius $p$ non-residuum. Si vero $p$ formae $4n+3$, erit $-A$
ipsius $p$ non-residuum.

Nam si ponamus $A = a^{\alpha} b^{\beta} c^{\gamma} \cdots$, ubi $a$, $b$, $c$ sunt numeri primi
positivi impares, positisque signis pro his numeris ut in art.~133, erit $-A$ ipsius
$p$ residuum vel non-residuum, prout $p$ ipsius $a^{\alpha} b^{\beta} c^{\gamma} \cdots$ est residuum vel
non-residuum. At quum $p$ sit ipsius $A$ non-residuum, saltem unus factor ipsius
$A$ erit, cuius $p$ non-residuum est, adeoque etiam $-A$ ipsius $p$ non-residuum.

\textbf{136.}
Theorema fundamentale pro numeris parvis verum esse, per inductionem fa-
cile confirmari, atque sielimes determinari potest usque ad quem certo locum teneat.
Hanc inductionem institutamesse postulamus: prorsus autem indifferens est quous-
que eam persequuti simus; sufficeret adeo, si tantummodo usque ad numerum $5$
eam confirmavissemus, hoc autem per unicam observationem absolvitur, quod est
$+5 \N\, 3$, $+3 \N\, 5$.

Iam si theorema fundamentale generaliter verum non est, dabitur limes ali-
quis, $T$, usque ad quem valebit, ita tamen ut usque ad numerum proxime maio-
rem, $T+1$, non amplius valeat. Hoc autem idem est ac si dicamus, dari duos
numeros primos, quorum maior sit $T+1$, et qui inter se comparati theoremati
fundamentali repugnent, binos autem alios numeros primos quoscunque, si modo
ambo ipso $T+1$ sint minores, huic theoremati esse consentaneos. Unde sequi-
tur, propositiones artt.~131, 132, 133 usque ad $T$ etiam locum habituras. Hanc
vero suppositionem consistere non posse nunc ostendemus. Erunt autem secun-
dum formas diversas quas tum $T+1$, tum numerus primus ipso $T+1$ minor,
quem cum $T+1$ comparatum theoremati repugnare supposuimus, habere pos-
sunt, casus sequentes distinguendi. Numerum istum primum per $p$ designamus.

Quando tum $T+1$ tum $p$ sunt formae $4n+1$, theorema fundamentale
duobus modis falsum esse posset, scilicet si simul esset,
vel $\pm p \R\, (T+1)$ et $\pm(T+1) \N\, p$
vel simul $\pm p \N\, (T+1)$ et $\pm(T+1) \R\, p$

Quando tum $T+1$ tum $p$ sunt formae $4n+3$, theor.\ fund.\ falsum erit,
si simul fuerit
vel $+p \R\, (T+1)$ et $-(T+1) \N\, p$
(sive quod eodem redit $-p \N\, (T+1)$ et $+(T+1) \R\, p$)
vel $+p \N\, (T+1)$ et $-(T+1) \R\, p$
(sive $-p \R\, (T+1)$ et $+(T+1) \N\, p$)

Quando $T+1$ est formae $4n+1$, $p$ vero formae $4n+3$, theor. fund. falsum erit, si fuerit
vel $\pm p \R\, (T+1)$ et $-(T+1) \N\,p$ (sive $-(T+1) \R\, p$)
vel $\pm p \N\, (T+1)$ et $-(T+1) \N\,p$ (sive $+(T+1) \R\, p$)


Quando $T+1$ est formae $4n+3$, $p$ vero formae $4n+1$, theor. fund. falsum erit, si fuerit
vel $\pm p \R\, (T+1)$ et $\pm(T+1) \N\, p$ (sive $-(T+1) \N\,p$)
vel $\pm p \N\, (T+1)$ et $\pm(T+1) \R\, p$ (sive $+(T+1) \R\,p$)

\textbf{137.}
Quando $T+1$ est formae $4n+1$, $p$ formae $4n+1$, atque $+p \N\, (T+1)$
sive $p \R\, (T+1)$, non poterit esse $+(T+1) \R\, p$ (Casus primus supra).

\textit{Demonstratio.} Sint omnia numerorum primorum ipso $T+1$ minorum
multitudines ut in art.~131; sitque $e$ par, per quem $p$ est non-residuum, atque
$r$, $r'$, $r''$ etc., $n$, $n'$, $n''$ etc.\ ut in art.~131. Quoniam $p$ est ipsius $T+1$
non-residuum, necesse est ut sit ipsius factorum alicuius ipsius $e$ non-residuum.
Sit $q$ factor primus ipsius $e$, cuius $p$ non-residuum est; eritque $q$ formae
$4n+3$ vel $4n+1$, prout $e$ est formae $4n+3$ vel $4n+1$. At per hyp.\ est
$p \R\, q$, adeoque per theor.\ fund.\ (quum $q < T+1$) erit $\pm q \R\, p$, accepto
signo prout $q$ formae $4n+1$ vel $4n+3$. Quare fiet $\pm e \R\, p$, signis perinde
compositis. Sit $e \equiv f^2 \pmod{p}$, ita ut $f$ sit impar et $< p$. Tum si esset
$+(T+1) \R\, p$, foret etiam $+(T+1)f^2 \R\, p$, adeoque $+e(T+1) \R\, p$.
At $e(T+1) = R$ est numerus ipso $T+1$ minor, qui per hyp.\ theorema
fundamentale confirmat. Quare etiam $p \R\, R$, adeoque per art.~105, $p \R\, e(T+1)$,
contra hyp.\ Igitur $+(T+1) \N\, p$.

\textbf{138.}
Quando $T+1$ est formae $4n+1$, $p$ formae $4n+1$, atque $+p \R\, (T+1)$,
non poterit esse $+(T+1) \N\, p$ (Casus secundus).

\textit{Demonstratio} prorsus similis praecedenti.

\textbf{139.}
Quando $T+1$ est formae $4n+1$, $p$ formae $4n+3$, atque $+p \N\, (T+1)$
sive $-p \R\, (T+1)$, non poterit esse $-(T+1) \N\,p$, i.~e.\ $-(T+1) \R\, p$
(Casus tertius).

\textit{Demonstratio.} Sit $e$ par, per quem $p$ est residuum, atque $q$ factor primus
ipsius $e$, cuius $p$ residuum est; eritque $q$ formae $4n+1$ vel $4n+3$. At per
hyp.\ est $p \R\, q$, adeoque per theor.\ fund.\ erit $\pm q \R\, p$. Hinc sicuti in
casu praecedente concluditur, esse $\pm e \R\, p$. Sit $e \equiv f^2 \pmod{p}$, $f$ impar.
Si esset $-(T+1) \R\, p$, foret $-e(T+1) \R\, p$, adeoque $-e(T+1)$ esset
residuum ipsius $p$, cuius $e$ est residuum; quare $-(T+1)$ esset residuum
ipsius $p$, contra hyp.

\textbf{140.}
Quando $T+1$ est formae $4n+1$, $p$ formae $4n+3$, atque $+p \R\, (T+1)$,
non poterit esse $-(T+1) \R\, p$ (Casus quartus).

\textit{Demonstratio} similis.

\textbf{141.}
Quando $T+1$ est formae $4n+3$, $p$ formae $4n+1$, atque $+p \N\, (T+1)$,
non poterit esse $+(T+1) \R\, p$ (Casus quintus).

\textit{Demonstratio.} Sit $e$ par, per quem $p$ est non-residuum; eritque $e$ formae
$4n+3$ vel $4n+1$, adeoque $-e$ formae $4n+1$ vel $4n+3$. Sit $q$ factor
ipsius $e$, cuius $p$ non-residuum; erit $q$ formae $4n+3$ vel $4n+1$. At per
hyp.\ $p \R\, q$, adeoque $\pm q \R\, p$. Hinc $\pm e \R\, p$, signis composite.
Si esset $+(T+1) \R\, p$, foret $-e(T+1) \R\, p$ (quoniam $-e$ formae $4n+1$);
at $-e(T+1) < T+1$, adeoque per hyp.\ theorema fund.\ confirmat, eritque
$p \R\, -e(T+1)$, adeoque $p \R\, (T+1)$, contra hyp.

\textbf{142.}
Quando $T+1$ est formae $4n+3$, $p$ formae $4n+1$, atque $+p \R\, (T+1)$,
non poterit esse $+(T+1) \N\, p$ (Casus sextus).

\textit{Demonstratio} similis.

\textbf{143.}
Quando $T+1$ est formae $4n+3$ ($= b$), $p$ formae $4n+3$, atque
$+p \N\, b$ sive $-p \R\, b$, non poterit esse $-b \N\, p$ (sive $+b \R\, p$)
(Casus septimus).

\textit{Demonstratio.} Sit $e$ par, per quem $p$ est residuum; eritque $e$ formae
$4n+1$ (quia alioquin $-e$ esset formae $4n+3$, adeoque $-ep \R\, b$, et per
hyp.\ cum $e < T+1$, theorema fund.\ confirmaret $-b \R\, e$, adeoque $-b \R\, p$,
contra id quod supponitur). Ergo $e$ formae $4n+1$. Sit $q$ factor ipsius $e$
cuius $p$ residuum; erit $q$ formae $4n+1$. At $p \R\, q$, adeoque $+q \R\, p$,
adeoque $+e \R\, p$. Si esset $+b \R\, p$, foret $+be \R\, p$, adeoque $be < T+1$
esset talis, cuius $p$ residuum, contra hyp.

\textbf{144.}
Casus octavus. Quando $T+1$ est formae $4n+3$ ($= b$), $p$ formae $4n+1$,
atque $+p \N\, b$ sive $p \R\, b$, non poterit esse $+b \R\, p$. (Casus ultimus supra).

Demonstratio perinde procedit ut in casu praecedenti.

\textbf{145.}
In demonstratt.\ praecc.\ semper pro $e$ valorem parem accepimus (artt.~137
--144); observare convenit, etiam valorem imparem adhiberi potuisse, sed tum plu-
res adhuc distinctiones introducendae fuissent. Qui bis disquisitionibus delec-
tantur haud inutile facient si vires suas in evolutione horum casuum exercitent.

Praeterea theoremata ad residua $+2$ et $-2$ pertinentia tunc supponi debuis-
sent; quum vero nostra demonstratio absque his theorematibus sit perfecta, no-
vam hinc methodum nanciscimur, illa demonstrandi. Quae minime est contem-
nenda, quum methodi, quibus supra pro demonstratione theorema $+2$ esse
residuum cuiusvis numeri primi formae $8n+1$, usi sumus, minus directae videri
possint. Reliquos casus (qui ad numeros primos formarum $8n+3$, $8n+5$,
$8n+7$ spectant) per methodos supra traditas demonstratos, illudque theorema
tantummodo per inductionem inventum esse supponemus; hanc autem inductio-
nem per sequentes reflexiones ad certitudinis gradum evehemus.

Si $+2$ omnium numerorum primorum formae $8n+1$ residuum non es-
set, ponatur minimus primus huius formae, cuius non-residuum $+2$, $= a$, ita
ut pro omnibus primis ipso $a$ minoribus theorema valeat. Tum accipiatur nu-
merus aliquis primus $< 2\sqrt{a}$, cuius non-residuum $a$ (qualem dari ex art.~129 fa-
cile deducitur). Sit hic $= p$, eritque per theor.\ fund.\ $p \N\, a$. Hinc fit $+2p \R\, a$.
Sit itaque $e^2 \equiv 2p \pmod{a}$, ita ut $e$ sit impar atque $< a$. Tum duo casus erunt
distinguendi.

I. Quando $e$ per $p$ non est divisibilis. Sit $e^2 = 2p + aq$; eritque $q$ po-
sitivus, formae $8n+7$ vel formae $8n+3$ (prout $p$ est formae $4n+1$ vel
$4n+3$), $< a$, atque per $p$ non divisibilis. Iam omnes factores primi ipsius $q$
in quatuor classes distribuantur, sint scilicet $e$ formae $8n+1$, $f$ formae $8n+3$,
$g$ formae $8n+5$, $h$ formae $8n+7$; productum e factoribus primae classis sit
$E$, producta e factoribus secundae, tertiae, quartae classis respective, $F$, $G$, $H$\footnote{Si ex aliqua classe nulli factores adessent, loco producti ex his $1$ scribere oporteret.}.
Iam patet, ex $q$ fieri $E \cdot F \cdot G \cdot H$. Sit $R$ numerus e factoribus ipsius $q$,
quorum $a$ sit non-residuum, eritque $R$ residuum ipsius $a$, adeoque etiam $qR$
residuum ipsius $a$, i.~e.\ $2pR$ residuum ipsius $a$. At per art.~131 foret
$+2 \R\, R$ (quia $R < a$), adeoque $pR \R\, a$, et $p \R\, a$, contra hyp.

II. Quando $e$ per $p$ est divisibilis. Sit $e = gp$, adeoque $g^2 p = 2 + \frac{aq}{p}$.
Manifesto $q$ impar et formae $8n+7$ vel $8n+3$, ut supra. Demonstratio
perinde procedit.

\textbf{146.}
Iam demonstrato theoremate fundamentali, sequitur ut theoremata de
residuis $+2$ et $-2$ etiam rigorose demonstrata sint. Hinc facile deducitur,
formam generalem divisorum et non-divisorum expressionis $x^2 - A$, de qua in
sequentibus agemus.

\subsection*{Formae divisorum ipsius $xx - A$.}
\addcontentsline{toc}{subsection}{Formae divisorum ipsius $xx - A$.}

\textbf{147.}
Omnis numerus, qui non dividitur per numerum primum datum $p$,
aut residuum aut non-residuum ipsius $p$ erit; et si $A$ est non-residuum ipsius
$p$, nullus numerus formae $x^2 - A$ per $p$ divisibilis erit, i.~e.\ $p$ erit in forma
non-divisorum expressionis $x^2 - A$. Si vero $A$ est residuum ipsius $p$, omnes
numeri formae $x^2 - A$ per $p$ divisibiles erunt, i.~e.\ $p$ erit in forma divisorum.
Quare si $A$ est productum e factoribus $a$, $b$, $c$ etc., numerusque primus $p$
ad $A$ primus, atque $r$, $r'$, $r''$ etc.\ designent omnes numeros, qui sunt non-residua
nullius numerorum $a$, $b$, $c$ etc.; $n$, $n'$, $n''$ etc.\ omnes qui sunt non-residua
unius e factoribus $a$, $b$, $c$ etc.; denique numeri, qui sunt non-residua duorum,
trium etc.\ eorundem factorum: tum omnes numeri primi in aliqua formarum
$4Ak + r$, $4Ak + r'$, $4Ak + r''$ etc.\ contenti erunt divisores expressionis
$x^2 - A$, omnes vero in aliqua formarum $4Ak + n$, $4Ak + n'$, $4Ak + n''$ etc. contenti erunt non-divisores eiusdem expressionis. Et vice versa, omnis divisor
aut non-divisor positivus atque alicuius ex numeris $a$, $b$, $c$ etc.\ residuum vel non-
residuum, hic ipse numerus ipsius $p$ residuum vel non-residuum erit (theor.\ fund.).
Quare si inter numeros $a$, $b$, $c$ etc.\ sunt $m$, quorum non-residuum est $p$, totidem
erunt non-residua ipsius $p$, adeoque si $p$ in aliqua formarum priorum continetur,
erit $m$ par et $A \R\, p$, si vero in aliqua posteriorum, erit $m$ impar atque $A \N\, p$.

\textit{Ex.} Sit $A = +105 = 3 \cdot 5 \cdot 7$. Tum numeri $r$, $r'$, $r''$ etc.\ erunt
hi: $1$, $4$, $16$, $46$, $64$, $79$ (qui sunt non-residua nullius numerorum $3$, $5$, $7$); $2$, $8$, $23$,
$32$, $53$, $92$ (qui sunt non-residua numerorum $3$, $5$); $26$, $41$, $59$, $89$, $101$, $104$ (qui sunt
non-residua numerorum $3$, $7$); $13$, $52$, $73$, $82$, $97$, $103$ (qui sunt non-residua nume-
rorum $5$, $7$). Numeri autem $n$, $n'$, $n''$ etc.\ erunt hi: $11$, $29$, $44$, $71$, $74$, $86$; $22$,
$37$, $43$, $58$, $67$, $88$; $19$, $31$, $34$, $61$, $76$, $94$; $17$, $38$, $47$, $62$, $68$, $83$. Seni primi sunt non-
residua ipsius $3$, seni posteriores non-residua ipsius $5$, tum sequuntur non-residua
ipsius $7$, tandem ii qui sunt non-residua omnium trium simul.

Facile ex combinationum theoria atque artt.~32, 96 deducitur, numerorum
$r$, $r'$, $r''$ etc.\ multitudinem fore
\[
\frac{1}{2}(1 + \alpha + \beta + \gamma + \alpha\beta + \alpha\gamma + \beta\gamma + \alpha\beta\gamma + \cdots)
\]
numerorum $n$, $n'$, $n''$ etc.\ multitudinem
\[
\frac{1}{2}(1 - \alpha + \beta - \gamma - \alpha\beta + \alpha\gamma - \beta\gamma + \alpha\beta\gamma - \cdots)
\]
ubi $l$ designat multitudinem numerorum $a$, $b$, $c$ etc.\ et utraque series continuanda donec abrumpatur. (Dabuntur scilicet $t$ numeri
qui sunt residua omnium $a$, $b$, $c$ etc., qui sunt non-residua duorum, etc. sed demonstrationem hanc fusius explicare brevitas non permittit). Utruisque
autem seriei summa\footnote{Neglecto factore $t$.} est $2^{l-1}$. Scilicet prior prodit ex hac
iungendo terminum secundum et tertium, quartum et quintum etc., posterior vero
ex eadem iungendo terminum primum atque secundum, tertium et quartum etc.

Dabuntur itaque tot formae divisorum ipsius $xx - A$, quot dantur formae non-
divisorum, scilicet $2^{l-1}(a-1)(b-1)(c-1)$ etc.

\textbf{148.}
Si adhuc requiritur ut omnes formae divisorum inter se sint diversae,
idem semper obtinebitur, si pro $a$, $b$, $c$ etc.\ accipiantur omnia numeri $A$ divisores
primi positivi impares, quorum multitudo $= l$, et signa pro his numeris prout
in art.~133 determinata adhibeantur. Tum enim facile perspicitur, formas di-
visorum productas e formis ipsorum $a$, $b$, $c$ etc.\ necessario diversas esse debere.
Sint enim $P$, $Q$ producta e factoribus diversis e numeris $a$, $b$, $c$ etc., atque
ponatur $P$ constare e factoribus $a^{\alpha} b^{\beta} c^{\gamma} \cdots$, $Q$ vero e factoribus
$a^{\alpha'} b^{\beta'} c^{\gamma'} \cdots$, ubi exponentes $\alpha$, $\beta$, $\gamma$, $\ldots$, $\alpha'$, $\beta'$, $\gamma'$, $\ldots$ sunt vel $0$ vel $1$,
sed non omnes inter se pares. Tum si formae divisorum pro $P$ et $Q$ essent eaedem,
foret $P \equiv Q \pmod{4A}$, adeoque $P - Q$ per $4A$ divisibilis. At facile perspi-
citur per continuationem huius operationis tandem ad numeros perventum iri,
qui diversas formas habere debent.

\textbf{149.}
Si $A$ est numerus primus formae $4n+1$, forma divisorum erit
$4Ak + r$, $4Ak + r'$, $4Ak + r''$ etc., ubi $r$, $r'$, $r''$ etc.\ sunt residua ipsius
$A$, et forma non-divisorum erit $4Ak + n$, $4Ak + n'$, $4Ak + n''$ etc., ubi
$n$, $n'$, $n''$ etc.\ sunt non-residua ipsius $A$. At si $A$ est formae $4n+3$,
forma divisorum ipsius $xx + A$ erit $4Ak + r$, ubi $r$ est residuum ipsius $A$,
forma non-divisorum erit $4Ak + n$, ubi $n$ est non-residuum ipsius $A$. Denique
si $A$ est formae $4n$, facile perspicitur, formas divisorum et non-divisorum per
$4A$ pendere a forma ipsius $A$ et a divisoribus ipsius $A$.

\textbf{150.}
Si habentur omnes formae divisorum et non-divisorum ipsius
$xx - A$ pro modulo $4A$, et quaeritur num numerus datus $N$ ad aliquam harum
formarum referendus sit, id quod vocatur indicium per quod $N$ ad aliquam
formarum refertur, facillime obtinetur ope theorematis fundamentalis. Sit enim
$N = pqrs\cdots$, ubi $p$, $q$, $r$, $s$ etc.\ sunt numeri primi (positivi) diversi, atque
signum numeri $N$ ita accipiatur ut sit $+N$, si $N$ formae $4n+1$, et $-N$,
si $N$ formae $4n+3$. Tum erit $N$ ipsius $A$ residuum, si omnes factores
$p$, $q$, $r$, $s$ etc.\ singuli per se respiciantur, secundum theor.\ fund.\ productum
omnium residuorum praebeant; non-residuum vero, si numerus non-residuorum
inter factors $p$, $q$, $r$, $s$ etc.\ impar sit.

\textbf{151.}
Sed haec de formis divisorum ipsius $xx - A$ hactenus. Iam ad alia
transeamus. Si $A$ est productum e factoribus primis $a$, $b$, $c$ etc., omnes nu-
meri primi, qui sunt formae $4Ak + r$, $4Ak + r'$ etc., erunt divisores ipsius
$xx - A$; omnes vero qui sunt formae $4Ak + n$, $4Ak + n'$ etc., erunt non-
divisores. Et vice versa, si $p$ est divisor aliquis ipsius $xx - A$, necesse est ut
$p$ sit in aliqua formarum $4Ak + r$, $4Ak + r'$ etc.\ contentus; si vero $p$ est
non-divisor, in aliqua formarum $4Ak + n$, $4Ak + n'$ etc.\ contentus erit.

\textbf{152.}
De congruentiis secundi gradus \textit{non puris}, i.~e.\ in quibus terminus
secundus non evanescit, hic pauca adiicimus. Sit proposita congruentia
\[
x^2 + 2bx + c \equiv 0 \pmod{m}.
\]
Haec aequipollet duabus congruentiis
\[
2x + 2b \equiv 0 \pmod{m} \quad \text{et} \quad x^2 + c - b^2 \equiv 0 \pmod{m},
\]
quarum prior est pura et methodis huius sectionis tractabilis. Si priori satisfit
per $x \equiv \alpha \pmod{m}$, substituendo in posteriori habetur $\alpha^2 + c - b^2 \equiv 0$,
adeoque $c - b^2$ debet esse residuum ipsius $m$, et quidem si $m$ est compositus,
residuum omnium factorum ipsius $m$.

Si vero congruentia non pura ita proponitur
\[
ax^2 + bx + c \equiv 0 \pmod{m},
\]
um per art.~106 reducitur ad formam puram, multiplicando per $4a$ et adiiciendo
$bb$:
\[
(2ax + b)^2 \equiv b^2 - 4ac \pmod{4am},
\]
quae si ponendo $2ax + b = y$ transformatur in
\[
y^2 \equiv b^2 - 4ac \pmod{4am},
\]
methodis supra explicatis solvenda est.


\section*{Sectio Quinta.}
\addcontentsline{toc}{section}{Sectio Quinta.}
\begin{center}
\Large\textbf{De Formis Aequationibusque Indeterminatis}\\
\Large\textbf{Secundi Gradus.}
\end{center}

\subsection*{Disquisitionis propositum: formarum definitio et signum.}
\addcontentsline{toc}{subsection}{Disquisitionis propositum.}

\textbf{153.}
In hac sectione imprimis de functionibus duarum indeterminatarum $x$, $y$,
huius formae
\[
axx + 2hxy + cyy
\]
ubi $a$, $h$, $c$ sunt integri dati, tractabimus, quas \textit{formas secundi gradus} sive simpli-
citer \textit{formas} dicemus. Huic disquisitioni superstruetur solutio problematis famosi,
invenire omnes solutiones aequationis cuiuscunque indeterminatae secundi gradus
duas incognitas implicantis sive hae incognitae valores integros sive rationales
tantum nancisci debeant. Problema hoc quidem iam ab ill.\ La Grange in omni
generalitate est solutum, multaque insuper ad naturam formarum pertinentia tum
ab hoc ipso magno geometra tum ab ill.\ Euler partim primum inventa, partim,
a Fermatio olim inventa, demonstrationibus munita. Sed nobis acriter formarum
perquisitioni insistentibus tam multa nova se obtulerunt, ut totum argumentum
ab integro resumere operae pretium duxerimus, eo magis, quod Virorum illorum
inventa, multis locis sparsa, paucis innotuisse experti sumus; porro quod metho-
dus, per quam haec tractabimus, nobis ad maximam partem est propria; tandem
quod nostra sine nova illorum expositione ne intelligi quidem possent. Nullum
tamen arbitramur alienum fore, si quae a viris illis iam inventa sunt, quam bre-
vissime attigerimus, atque eo ordine, qui nobis in sequentibus maxime utilis vide-
tur.

\textbf{154.}
Si $a$, $h$, $c$ sunt numeri integri, formam $axx + 2hxy + cyy$ per $f$
designabimus. Formae datae $f$ \textit{determinans} vocabimus numerum $D = h^2 - ac$,
quem etiam \textit{determinantem} formarum $a$, $2h$, $c$ dicemus. Duo numeri, e quibus
determinans componitur, $a$, $2h$, $c$, vocabuntur \textit{coefficientes} formae; $a$ et $c$ etiam
\textit{termini extremi}, $2h$ \textit{terminus medius} vocabitur.

Si $D$ est negativus, omnes formae, quarum determinans est $D$, vocabimus
\textit{formas definitas}; forma autem $f$ dicetur \textit{positiva definita} si $a$ et $c$ sunt ambo
positivi, \textit{negativa definita} si $a$ et $c$ sunt ambo negativi. Si $D$ est positivus,
vocabimus \textit{formas indefinitas}; forma autem $f$ dicetur \textit{indefinita} si $a$ et $c$ sunt
diversorum signorum. Si $D = 0$, dicetur forma \textit{forma anceps}.

Si forma $f$ transformatur in formam $f'$ per substitutionem
\[
x = \alpha x' + \beta y', \qquad y = \gamma x' + \delta y'
\]
umbi $\alpha$, $\beta$, $\gamma$, $\delta$ sunt integri, forma $f'$ dicitur \textit{contenta} sub forma $f$, et forma
$f$ dicitur \textit{continere} formam $f'$. Si simul forma $f'$ continet formam $f$, formae
$f$ et $f'$ dicuntur \textit{aequivalentes}. Si autem forma $f$ continet formam $f'$ sed $f'$
non continet formam $f$, dicitur forma $f$ \textit{implicare} formam $f'$, et forma $f'$ \textit{implicari}
a forma $f$.

\textbf{155.}
Si forma $f$ transit in formam $f'$ ponendo
\[
x = \alpha x' + \beta y', \qquad y = \gamma x' + \delta y'
\]
erit determinans formae $f'$ aequalis determinanti formae $f$ multiplicato per
$(\alpha\delta - \beta\gamma)^2$, adeoque si forma $f'$ sub forma $f$ contenta est, determinans
ipsius $f'$ per determinans ipsius $f$ divisibile erit. Si formae $f$ et $f'$ sunt aequi-
valentes, determinantes erunt aequales\footnote{Manifestum est ex analysi praecedente hanc propositionem etiam ad formas, quarum determinans $= 0$, patere. Sed aequatio $(\alpha\delta - \beta\gamma)^2 = 1$ ad hunc casum non est extendenda.} atque $(\alpha\delta - \beta\gamma)^2 = 1$. In hoc casu
formas aequivalentes dicemus. Quare ad formarum aequivalentiam aequalitas de-
terminantium est conditio necessaria, licet illa ex hac sola minime sequatur.

Substitutionem $x = \alpha x' + \beta y'$, $y = \gamma x' + \delta y'$ vocabimus \textit{transformationem pro}-
\textit{priam}, si $\alpha\delta - \beta\gamma$ est numerus positivus, \textit{impropriam}, si $\alpha\delta - \beta\gamma$ est negativus;
formam $f'$ \textit{proprie} aut \textit{improprie} sub forma $f$ contentam esse dicemus, si
per transformationem propriam aut impropriam in formam $f'$ transmutari potest.

Si itaque formae $f$, $f'$ sunt aequivalentes, erit $(\alpha\delta - \beta\gamma)^2 = 1$, adeoque si
transformatio est propria, $\alpha\delta - \beta\gamma = +1$, si est impropria, $\alpha\delta - \beta\gamma = -1$. Si
plures transformationes simul sunt propriae, aut simul impropriae, \textit{similes} eas di-
cemus; propriam contra et impropriam \textit{dissimiles}.

\subsection*{Aequivalentia propria et impropria.}
\addcontentsline{toc}{subsection}{Aequivalentia propria et impropria.}

\textbf{156.}
Si formarum $F$, $F'$ determinantes sunt aequales atque $F'$ sub $F$ contenta: etiam
$F$ sub $F'$ contenta erit et quidem proprie vel improprie, prout $F'$ sub $F$ proprie vel
improprie continetur. Transeat $F$ in $F'$ ponendo
\[
x = \alpha x' + \beta y', \qquad y = \gamma x' + \delta y'
\]
transibitque $F'$ in $F$ ponendo
\[
x' = \delta x - \beta y, \qquad y' = -\gamma x + \alpha y
\]
Patet enim per hanc substitutionem ex $F'$ fieri idem quod fiat ex $F$ ponendo,
\[
x = \alpha(\delta x - \beta y) + \beta(-\gamma x + \alpha y), \quad y = \gamma(\delta x - \beta y) + \delta(-\gamma x + \alpha y)
\]
sive
\[
x = (\alpha\delta - \beta\gamma)x, \qquad y = (\alpha\delta - \beta\gamma)y
\]
Hinc vero manifeste ex $F$ fit $(\alpha\delta - \beta\gamma)^2 F$, i.~e.\ rursus $F$ (art.\ praec.). Perspi-
cuum autem est, transformationem posteriorem esse propriam vel impropriam,
prout prior sit propria vel impropria.

Si tum $F'$ sub $F$, tum $F$ sub $F'$ proprie continetur, formas \textit{proprie aequi-}
\textit{valentes} vocabimus; si tum $F'$ sub $F$, tum $F$ sub $F'$ improprie continetur, \textit{improprie}
\textit{aequivalentes}; si alterutra continentia propria, altera impropria est, \textit{semi-aequivalentes}.

\textbf{157.}
Si formae $F$, $F'$ eundem determinantem habent, atque $F'$ sub $F$ conti-
nentur, forma $F'$ transibit in formam $G$, quae etiam sub $F$ continebitur, si
ponamus
\[
x = mX + tY, \qquad y = nX + uY
\]
umbi $m$, $n$, $t$, $u$ sunt integri quicunque. Tum si $F$ transit in $F'$ ponendo
$x = \alpha x' + \beta y'$, $y = \gamma x' + \delta y'$, erit $F'$ transmutata in $G$ ponendo
\[
x' = mX + tY, \qquad y' = nX + uY.
\]
Quare $G$ sub $F'$ contenta erit, adeoque etiam sub $F$. Manifesto si $F$ et $F'$
valentes, si illae sub invicem improprie, vocabimus improprie aequivalentes. Ce-
terum formae proprie aequivalentes semper etiam improprie aequivalentes erunt,
si modo forma anceps detur, cuius determinans sit aequalis determinanti formae
propositae.

\subsection*{Formae contiguae.}
\addcontentsline{toc}{subsection}{Formae contiguae.}

\textbf{158.}
Si forma $axx + 2hxy + cyy$ ita transformata est, ut coefficientem
ipsius $xx$ servet, coefficientem vero ipsius $xy$ mutet, dicitur forma \textit{contigua} priori.
Sit forma $axx + 2hxy + cyy$ transformata in $a'x'^2 + 2h'x'y' + c'y'^2$; si
$a' = a$, formae dicuntur \textit{contiguae}. Si denique forma $f$ transit in formam $f'$
per substitutionem $x = x'$, $y = y' + mx'$, formae $f$ et $f'$ dicuntur \textit{contiguae},
quoniam coefficientem ipsius $xx$ servant. Hinc sequitur, maximum divisorem communem numerorum $a$, $2h$, $c$ si-
per $d$, $2h'$, $c'$ designamus, ex aequ.\ praecc.\ sequentes novas deducemus\footnote{In his formulis, si pro $a$, $h$, $c$ valores ex $1$, $3$, $5$ substituuntur, fit}:
\begin{align*}
a' &= a, \\
h' &= h + am, \\
c' &= am^2 + 2hm + c.
\end{align*}
Hinc patet, $a$ esse divisorem communem numerorum $a$, $2h'$, $c'$; quare si $a$ ad
$2h$ primus est, erit etiam $a'$ ad $2h'$ primus. Porro facile perspicitur, determi-
nantem formae $f'$ aequalem esse determinanti formae $f$, uti ex natura aequiva-
lentiae sequitur.

\textbf{159.}
Si forma $f$ transit in formam $f'$ per transformationem propriam, atque
$f'$ in formam $f''$ per transformationem propriam, tum etiam $f$ transibit in $f''$
per transformationem propriam. Nam si $f$ transit in $f'$ ponendo
$x = \alpha x' + \beta y'$, $y = \gamma x' + \delta y'$, atque $f'$ in $f''$ ponendo
$x' = \alpha' x'' + \beta' y''$, $y' = \gamma' x'' + \delta' y''$,
tum $f$ transit in $f''$ ponendo
\begin{align*}
x &= (\alpha\alpha' + \beta\gamma')x'' + (\alpha\beta' + \beta\delta')y'', \\
y &= (\gamma\alpha' + \delta\gamma')x'' + (\gamma\beta' + \delta\delta')y''.
\end{align*}
Manifesto huius transformationis determinans erit
\[
(\alpha\alpha' + \beta\gamma')(\gamma\beta' + \delta\delta') - (\alpha\beta' + \beta\delta')(\gamma\alpha' + \delta\gamma')
= (\alpha\delta - \beta\gamma)(\alpha'\delta' - \beta'\gamma'),
\]
adeoque $= +1$, quoniam uterque factor per hyp.\ $= +1$. Q.~E.~D.

Simili modo, si forma $f$ transit in $f'$ per transformationem impropriam, et
$f'$ in $f''$ per transformationem impropriam, tum $f$ transit in $f''$ per transforma-
tionem \textit{propriam}; si vero alterutra est propria, altera impropria, tum $f$ transit
in $f''$ per transformationem impropriam.

\textbf{160.}
Si forma $f$ implicat formam $f'$, et haec implicat formam $f''$, tum etiam
$f$ implicat $f''$. Nam si $f'$ sub $f$ contenta est, et $f''$ sub $f'$ contenta est, manifesto
$f''$ etiam sub $f$ contenta erit. Quodsi $f'$ non sub $f$ contenta est, i.~e.\ $f'$
non transit in $f$ per substitutionem integram, manifesto $f''$ sub $f'$ contenta erit.

\subsection*{Numerorum repraesentatio.}
\addcontentsline{toc}{subsection}{Numerorum repraesentatio.}

\textbf{161.}
Dicitur numerus $M$ \textit{repraesentari} per formam $f$, si dentur numeri
integri $x$, $y$, quibus substitutis forma $f$ praebeat valorem $M$. Si $M$ repraesentatur
per formam $f$, et $f$ transit in formam $f'$ per substitutionem $x = \alpha x' + \beta y'$,
$y = \gamma x' + \delta y'$, manifesto $M$ etiam repraesentabitur per formam $f'$, si modo
$x'$, $y'$ valores integros habeant. Si $M$ repraesentatur per formam $f$ ponendo
$x = m$, $y = n$, ubi $m$, $n$ sunt integri ad se invicem primi, dicimus repraesenta-
tionem esse \textit{primam}; si vero $m$, $n$ divisorem communem habent, repraesentationem
esse \textit{secundariam}.

Si forma $axx + 2hxy + cyy$ repraesentat numerum $M$ ponendo $x = m$,
$y = n$, facile perspicitur esse $Mm = am^2 + 2hmn + cn^2$, adeoque
$D \cdot M = (hm + cn)^2 - a \cdot M$, si $D = h^2 - ac$.

\textbf{162.}
Si forma $f$ transit in formam $f'$ per substitutionem propriam, omnis
numeri $M$ repraesentatio per formam $f'$ eiusdem naturae est, ac repraesentatio
per formam $f$. Nam si $M = am^2 + 2hmn + cn^2$, ubi $m$, $n$ sunt integri,
tum ponendo $x = m$, $y = n$ in forma $f$, et $x' = \delta m - \beta n$,
$y' = -\gamma m + \alpha n$ in forma $f'$, habebitur eadem repraesentatio. Et vice versa,
omnis repraesentatio per $f'$ dat repraesentationem per $f$.

Si vero $f$ transit in $f'$ per substitutionem impropriam, omnis repraesentatio
per $f'$ dat repraesentationem per $f$, sed non vice versa; repraesentationes autem,
quae per $f$ dantur non vero per $f'$, dicuntur \textit{exclusae}.

\textbf{163.}
Formae \textit{ancipites.} Si forma $axx + 2hxy + cyy$ transit in formam
$a'x'^2 + 2h'x'y' + c'y'^2$ per substitutionem
\[
x = \alpha x' + \beta y', \qquad y = \gamma x' + \delta y',
\]
umbi $\alpha\delta - \beta\gamma = 0$, forma dicitur \textit{anceps}. In hoc casu forma anceps est, quando
$\alpha\delta = \beta\gamma$, adeoque forma non amplius duas indeterminatas habet, sed ad formam
simplicem reducitur.

Si formae $F$ et $G$ sunt ancipites, atque $F$ transit in $G$ per substitutionem
$x = mX + nY$, $y = tX + uY$, facile perspicitur fore
\begin{align*}
M &= am^2 + 2hmn + cn^2, \\
N &= amt + h(mu + nt) + cnu, \\
P &= at^2 + 2htu + cu^2.
\end{align*}
Porro per evolutionem facile confirmatur, esse
\begin{align*}
2e + 2a &= (\alpha - \alpha')(\delta + \delta') - (\beta - \beta')(\gamma + \gamma'), \\
2b &= (\alpha + \alpha')(\delta - \delta') - (\beta + \beta')(\gamma - \gamma').
\end{align*}
Hinc quoniam $m(\gamma + \gamma') = n(\alpha + \alpha')$, $m(\delta + \delta') = n(\beta + \beta')$ erit
$m(2e + 2a) = 2nb$ sive
\[
me + ma = nb \tag{7}
\]
Eodem modo erit
\[
ne = na + mc \tag{8}
\]
Iam si ad $mm(b\mu\mu - 2a\mu\nu - c\nu\nu)$ additur
\[
(1 - m\mu - n\nu)\{m\nu(e+a) + (m\mu + 1)b\}
\]
quod manifeste propter
\[
1 - m\mu - n\nu = 0, \quad me + ma = nb, \quad ne = na + mc
\]
prodit productis rite evolutis partibusque se destruentibus deletis, $2m\nu e + b$.
Quare erit
\[
mm(b\mu\mu - 2a\mu\nu - c\nu\nu) = 2m\nu e + b \tag{9}
\]
Eodem modo addendo ad $mn(b\mu\mu - 2a\mu\nu - c\nu\nu)$ haec:
\[
(1 - m\mu - n\nu)\{(m\nu + n\mu)e + b\}
\]
$+ (me + ma + nh)m\mu\nu + (ne - na + mc)n\nu\nu$
invenitur
\[
mn(b\mu\mu - 2a\mu\nu - c\nu\nu) = (m\nu + n\mu)e + a \tag{10}
\]
Denique addendo ad $nn(b\mu\mu - 2a\mu\nu - c\nu\nu)$ haec:
\[
(m\mu + n\nu - 1)\{(n\mu - m\nu)e + (n\nu + 1)b\}
\]
$+ (me + ma + nh)n\mu\nu + (ne - na + mc)(n\nu\nu + \nu)$
fit
\[
nn(b\mu\mu - 2a\mu\nu - c\nu\nu) = 2n\mu e + c \tag{11}
\]
Iam ex $9$, $10$, $11$, deducitur
\begin{align*}
&(Amm + 2Bmn + Cnn)(b\mu\mu - 2a\mu\nu - c\nu\nu) \\
&\quad = 2e(Am^2 + B(m\mu + n\nu) + Cn^2) + Ab + 2Ba + Cc
\end{align*}
sive propter
\[
[6] \qquad M = Amm + 2Bmn + Cnn, \quad Nr = (Am\nu + B(m\mu + n\nu) + Cn\mu)e
\]
\[
M(b\mu\mu - 2a\mu\nu - c\nu\nu) = 2Nr \tag{Q.E.D.}
\]

\textbf{164.}
Si forma $G$ est anceps, et transit in formam $F'$ per substitutionem
propriam, tum $F'$ etiam erit anceps. Quoniam $G$ transit in $F$ ponendo
$\xi = \mu x + \nu y$, $\eta = \rho x + \sigma y$,
forma $G$ per substitutionem $(S)$ transmutabitur in eandem formam, in quam $F$
transformatur ponendo
\begin{align*}
x &= \mu\{(\alpha\alpha' + \nu\gamma')x' + (\alpha\beta' + \nu\delta')y'\} + \nu\{(n\alpha' - m\gamma')x' + (n\beta' - m\delta')y'\} \\
&= \alpha(m\mu + n\nu)x' + \beta(m\mu + n\nu)y' \text{ sive } \alpha x' + \beta y'
\end{align*}
et
\begin{align*}
y &= \rho\{(\alpha\alpha' + \nu\gamma')x' + (\alpha\beta' + \nu\delta')y'\} + \sigma\{(n\alpha' - m\gamma')x' + (n\beta' - m\delta')y'\} \\
&= \gamma(m\mu + n\nu)x' + \delta(m\mu + n\nu)y' \text{ sive } \gamma x' + \delta y'
\end{align*}
Per hanc vero substitutionem $F$ transit in $F'$; quare per substitutionem $(S)$
$G$ etiam transibit in $F'$.

1.~Ex valoribus ipsorum $e$, $h$, $d$ invenitur $\alpha'e + \gamma\delta = 0$, sive propter
$d = -a$, $\alpha e + \alpha\alpha' + \beta\gamma = 0$; hinc ex $[7]$, $n\alpha e + n\alpha\alpha' = m\gamma e + m\gamma\alpha'$
sive
\[
(n\alpha - m\gamma)a = (m\gamma - n\alpha)e \tag{12}
\]
Porro fit $\alpha n b = \alpha m(e+a)$, $\gamma m b = m(\alpha e + \alpha\alpha')$ adeoque
\[
-(n\alpha - m\gamma)b = (\alpha - \alpha')me \tag{13}
\]
Denique fit $\gamma'e - \gamma a + \alpha c = 0$: hinc multiplicando per $n$, et pro $n\alpha$ substi-
tuendo valorem ex $[8]$ fit
\[
(n\alpha - m\gamma)c = (\gamma - \gamma')ne \tag{14}
\]
Simili modo eruitur $\beta'e + \delta b - \beta d = 0$, sive $n\beta'e + n\beta b + n\beta a = 0$, adeoque
per $[7]$, $n\beta'e + n\beta a = m\delta e + m\delta a$ sive
\[
(n\beta - m\delta)a = (m\delta - n\beta')e \tag{15}
\]
Porro fit $\beta n b = \beta m(e+a)$, $\delta m b = m(\beta'e + \beta a)$ adeoque
\[
(n\beta - m\delta)b = (\delta' - \delta)me \tag{16}
\]
Tandem $\delta'e - \delta a + \beta c = 0$: hinc multiplicando per $n$ et substituendo pro $n\alpha$
valorem ex $[8]$ fit
\[
(n\beta - m\delta)c = (\delta - \delta')ne \tag{17}
\]
Iam quum divisor communis maximus numerorum $a$, $b$, $c$ sit $r$, integri
$\mathfrak{A}$, $\mathfrak{B}$, $\mathfrak{C}$ ita accipi possunt, ut fiat
\[
\mathfrak{A} a + \mathfrak{B} b + \mathfrak{C} c = r
\]
Quo facto erit ex $12$, $13$, $14$; $15$, $16$, $17$
\begin{align*}
\mathfrak{A}(m\gamma - n\alpha) + \mathfrak{B}(\alpha - \alpha')m + \mathfrak{C}(\gamma - \gamma')n &= -(n\alpha - m\gamma), \\
\mathfrak{A}(n\alpha - m\gamma) + \mathfrak{B}(\delta' - \delta)m + \mathfrak{C}(\delta - \delta')n &= -(n\beta - m\delta)
\end{align*}
adeoque $-(n\alpha - m\gamma)$, $-(n\beta - m\delta)$ integri. Q.~E.~D.

\end{document}
